\section{Scope: a composite-source descent, not initial occupancy}
The preceding prime-shell trace theorem \cite{PrimeTrace} supplies an integer
return from an integral rational first trace when the distinguished denominator
is prime. It explicitly leaves the composite case open. Here the distinguished
denominator may be arbitrary. We classify every residual-divisor return that
keeps its original divisor word. Nine words have a uniform one-step return;
for every other fixed word, infinitely many hard primes give exact obstructions
within this transformation system. The original prime is never changed.

This remains the programme's Turn~7. No theorem below asserts that a prescribed
prime supplies the initial rational trace candidate. The Type-I/II formulas are
classical \cite[\S2, Propositions 2.2 and 2.6]{ET}; the trace relaxation is
inherited from \cite[TR17--19]{Trace}. The finite capacity bound used in the
optimality theorem is from \cite[Theorem ``Sharp finite-repair bound'']{Capacity}
and is reproved here. No historical-priority determination is made.
Here a \emph{hard prime} means
\[
 p\bmod840\in\{1,121,169,289,361,529\}.
\]
The counterfamily below lies in the first of these classes.

\section{The full original divisor box and its rational trace relaxation}
Fix a prime $p\equiv1\pmod4$, an integer $p/4<a<p/2$, and
$R=4a-p$. Thus $3\le R<p$, $R\equiv3\pmod4$ and $\gcd(pa,R)=1$.
Retain the integer word and every exponent:
\[
 a=\prod_l l^{e_l},\qquad u=\prod_l l^{f_l}\mid a^2,
 \quad 0\le f_l\le2e_l,\qquad \beta_l=f_l-e_l.
\]
For $c\in\{E,M\}$ define
\[
 G_E=4u+1,\qquad G_M=p+4u.
\]
Before requiring the full gate $R\mid G_c$, these original words give the
positive rational denominator triples
\begin{align}
 E:&\quad \left(a,\frac{pa+a^2/u}{R},\frac{pa+p^2u}{R}\right),\label{eq:Erational}\\
 M:&\quad \left(a,\frac{p(a+a^2/u)}{R},\frac{p(a+u)}{R}\right).\label{eq:Mrational}
\end{align}
They satisfy $4/p=1/x+1/y+1/z$ by multiplication.

\begin{lemma}[Exact trace domain]\label{lem:trace}
The tail trace $y+z$ is integral exactly when $R\mid G_c^2$.
The two reduced tail denominators then equal
\begin{equation}\label{eq:d}
 d=\frac R{\gcd(R,G_c)},\qquad d^2\mid R.
\end{equation}
They are integers exactly when $d=1$.
\end{lemma}
\begin{proof}
The traces are respectively $(a+pu)^2/(Ru)$ and $p(a+u)^2/(Ru)$.
Since $u\mid a^2$, both $u\mid(a+pu)^2$ and $u\mid(a+u)^2$.
All $p,a,u$ are units modulo $R$, in particular $\gcd(u,R)=1$.
Since $p\equiv4a\pmod R$,
$a+pu\equiv a(1+4u)$ and $4(a+u)\equiv p+4u$ there.
The exact tail denominators in \eqref{eq:Erational} and
\eqref{eq:Mrational} follow by multiplying their integer numerators by the
unit $u$ when needed.  If $l^e\parallel R$ and $g=v_l(G_c)$, the exponent
of $l$ in the reduced denominator is $\max(e-g,0)$.  The trace gate
$R\mid G_c^2$ gives $e\le2g$, hence
$2\max(e-g,0)\le e$.  This proves $d^2\mid R$ prime by prime.
\end{proof}
A record with $d>1$ will be called a proper trace source. It contains an
\emph{available integer divisor word}, but the full residual gate is not yet
satisfied. It is different from the earlier raw integer sum-gate source in
\cite[TR7--16]{Trace}, which imposed its full gate before testing a trace.

These sources also recover every rational integral-trace tail pair.  Let
$S=y+z\in\mathbb Z$ and put $b=Ry-pa$, $c=Rz-pa$.  The reciprocal identity
gives
\[
 yz=\frac{paS}{R},\qquad b+c=RS-2pa\in\mathbb Z,
 \qquad bc=(pa)^2\in\mathbb Z.
\]
Thus the rational numbers $b,c$ are roots of the monic polynomial
\[
 X^2-(RS-2pa)X+(pa)^2\in\mathbb Z[X].
\]
A rational algebraic integer is an integer.  Moreover
$1/y,1/z<R/(pa)$, so $b,c>0$.  The three possible pairs of $p$-valuations
are $(0,2),(1,1),(2,0)$.
For completeness, write the common reduced denominator as
$y=A/d$, $z=C/d$; it is common because $z=S-y$.  Then
$RAC=paSd^2$ and $\gcd(AC,d)=1$, so $d^2\mid R$.  If $R=d^2r$, then
$rAC=paS$; since $\gcd(r,pa)=1$, write $S=rT$.  The equality
$AC=paT$ and the facts $\gcd(AC,d)=\gcd(pa,d)=1$ give $\gcd(T,d)=1$.
Consequently $\gcd(R,S)=r$ and
\[
 d^2=R/\gcd(R,S).
\]
After retaining an E orientation, the first pair has
$b=a^2/u,c=p^2u$; the middle pair has $b=pa^2/u,c=pu$.
Their inverses are exactly \eqref{eq:Erational}--\eqref{eq:Mrational}.
Both middle tail orientations remain separate.

For any positive integer $n$ write
\[
 K(n)=\prod_{l^e\parallel n}l^{\lceil e/2\rceil},\qquad
 B(n)=\prod_{l^e\parallel n}l^{\lfloor e/2\rfloor}.
\]
The actual word budget is $u\mid a^2\iff K(u)\mid a$.
Put $m=4K(u)$. Then
\begin{equation}\label{eq:budget}
 p+R\equiv0\pmod m,\qquad \gcd(pR,m)=1.
\end{equation}
No absent prime coordinates are adjoined to this budget.

\section{Every same-word residual-divisor return}
\begin{theorem}[Complete divisor-removal domain]\label{thm:all}
For a trace source $(p,a,R,u,c)$, the complete set of original same-channel
returns that keep $p,u$ and replace $R$ by a divisor is indexed by the positive
integers
\begin{equation}\label{eq:domain}
 \boxed{\mathcal K=\{k\in\mathbb Z_{>0}:k\mid R,\ d\mid k,\
 k\equiv1\pmod m\}.}
\end{equation}
The map and the inverse after this source is retained are
\begin{equation}\label{eq:map}
 R'=R/k,\qquad a'=(p+R')/4,\qquad u'=u;
 \qquad k=R/R'.
\end{equation}
For a proper source every return has $a'<a$ and hence terminates in one
strictly decreasing step at an original integral state.
\end{theorem}
\begin{proof}
The full new gate $R/k\mid G_c$ is equivalent prime by prime to $d\mid k$.
The new integral distinguished denominator and its word budget together are
$p+R/k\equiv0\pmod m$. Multiplying by $k$ and using \eqref{eq:budget}
proves that this is equivalent to $k\equiv1\pmod m$.
This condition includes $k\equiv1\pmod4$, so $R'\equiv3\pmod4$ and $R'\ge3$.
Also $R'\le R<p$; hence $p/4<a'<p/2$. The condition $K(u)\mid a'$ proves
$u\mid a'^2$ in the actual integers. If $d>1$ then $k\ge d>1$ and $a'<a$.
All steps are equivalences, proving completeness.
\end{proof}

The full channel-labelled output is computed, not inherited unchanged from the
rational tails.  The label order records which tail recovers $u$; it is not a
claim that the denominators increase numerically.  Put
\[
 g'=\gcd(a',u),\quad h'=g'^2/u,\quad r'=u/g',\quad s'=a'/g'.
\]
Prime valuations prove $a'=h'r's'$, $u=h'r'^2$ and $\gcd(r',s')=1$.
For E, set $\kappa'=(pr'+s')/R'$ and return
$(a',h's'\kappa',ph'r'\kappa')$. For M, set
$\lambda'=(r'+s')/R'$ and return $(a',ph's'\lambda',ph'r'\lambda')$.
The full gates prove the quotient integrality. Their labelled divisor inverses
are $a'^2/(R'y-pa')$ and $pa'^2/(R'y-pa')$ respectively.
The centred prime-exponent vectors satisfy
\[
 \beta'_l=\beta_l+v_l(a)-v_l(a').
\]
A numerical word may survive even though the distinguished denominator's other
prime factors change. Those old and new prime coordinates remain recorded.

\subsection{Complete inverse fibres and integer coefficient maps}
Fix a full original target $(p,R',a',u,c)$ satisfying
$R'=4a'-p$, $3\le R'<p$, $R'\equiv3\pmod4$, $u\mid(a')^2$, and
$R'\mid G_c$.  Every proper input mapping to it is recovered by the positive
integers
\begin{equation}\label{eq:fibre}
 1\le k\le\left\lfloor\frac{p-1}{R'}\right\rfloor,
 \quad k\equiv1\pmod m,\quad R'k\mid G_c^2,
 \quad R'k\nmid G_c,
\end{equation}
with $R=R'k$ and $a=(p+R)/4$. The old word budget follows from $k\equiv1\pmod m$.
Indeed $p+R'\equiv0\pmod m$ at the target, while
$p+R'k=(p+R')+R'(k-1)\equiv0\pmod m$.  Thus $a$ is an integer and
$u\mid a^2$.  The bound and $k\equiv1\pmod4$ give $0<R<p$ and
$R\equiv3\pmod4$; the last two divisibilities in \eqref{eq:fibre} give the
trace gate and properness.  Finally $R'\mid G_c$ implies, prime by prime, that
the deletion denominator $d=R/\gcd(R,G_c)$ divides $k$.  These conditions are
therefore both necessary and sufficient, including all original exponent
multiplicities.

For exact fibre typing, let
\[
 \mathcal A=\{(s,k):s\text{ is a proper source and }k\in\mathcal K_s\},
 \qquad \mathcal T_{\rm hit}=F(\mathcal A),
\]
where $F(s,k)$ is the target in \eqref{eq:map}.  One source can have several
arrows and targets; for example $(p,a,R,u,c)=(101,41,63,1,M)$ has
$k=9$ and $k=21$.  The correctly typed pushforward is
\[
 F_*:\mathbb Z^{(\mathcal A)}\longrightarrow
 \mathbb Z^{(\mathcal T_{\rm hit})},\qquad [\alpha]\longmapsto[F(\alpha)].
\]
Choose one arrow $\alpha_t\in F^{-1}(t)$ for every hit target.  Then
\[
 \ker F_*=\bigoplus_{t\in\mathcal T_{\rm hit}}
 \left\langle[\alpha]-[\alpha_t]:
 \alpha\in F^{-1}(t)\setminus\{\alpha_t\}\right\rangle_{\mathbb Z}.
\]
The chosen arrows define a section $s([t])=[\alpha_t]$.  The actual integral
isomorphism retaining every arrow is
\[
 x\longmapsto\bigl(F_*x,\,x-sF_*x\bigr),\qquad
 \mathbb Z^{(\mathcal A)}\cong
 \mathbb Z^{(\mathcal T_{\rm hit})}\oplus\ker F_*,
\]
with inverse $(y,z)\mapsto s(y)+z$.  After the source and fibre-difference
coordinates are discarded, no inverse remains.

There is a finite coefficient calculation before any successful word is selected.
Use the group ring of $G_m=(\mathbb Z/m\mathbb Z)^\times$, with $[g]$ for its
basis elements.  The integer $d$ and every prime dividing $R/d$ are units
modulo $m$, because $\gcd(R,m)=1$.  Since $k=dv$ in \eqref{eq:domain},
the labelled pushforward is
\[
 \pi_{R,d}:\mathbb Z^{(\operatorname{Div}(R/d))}\longrightarrow
 \mathbb Z[G_m],\qquad e_v\longmapsto[v\bmod m].
\]
The group-ring image alone does not retain the divisor label $v$; the labelled
certificate does.  The complete integer count is
\begin{equation}\label{eq:poly}
 |\mathcal K|=\operatorname{coeff}_{[d^{-1}]}
 \prod_{l^e\parallel R/d}
       \left([1]+[l]+\cdots+[l^e]\right).
\end{equation}
Each factor is geometric: its $e+1$ exponent choices occur once each. It is not
an expansion of $([1]+[l])^e$ which would introduce coloured multiplicities.
The integer coefficient and the original divisors representing it are retained.

If a target word exists, one exists with
\begin{equation}\label{eq:length}
 \Omega(k/d)\le |G_m|-1=\varphi(m)-1.
\end{equation}
Choose a shortest original prime-occurrence word for $d^{-1}$. Its successive
partial products, starting with one, are distinct; a repeated partial product
would let an identity subword be removed without exceeding any exponent bound.
This proves \eqref{eq:length}. It bounds an existing word; it does not assert
that the target coefficient is nonzero.

\section{Square removal and the exact nine-seed guarantee}
Write
\[
 R=\delta t^2,\quad\delta=\prod_{l^e\parallel R}l^{e\bmod2},\qquad
 q_0=K(d),\qquad T=t/q_0.
\]
The squarefree part $\delta$ is not the radical, and it need not be coprime to $t$.
The fact $d^2\mid R$ proves $q_0\mid t$.
\begin{proposition}[Every square deletion and its capacity]\label{prop:square}
For $q\in\mathbb Z_{>0}$, the square members $k=q^2$ of
\eqref{eq:domain} are precisely
\begin{equation}\label{eq:square}
 q\mid t,\qquad q_0\mid q,\qquad q^2\equiv1\pmod m.
\end{equation}
Let $H_m=\{x^2:x\in G_m\}$. Their complete count is
\begin{equation}\label{eq:squarepoly}
 \operatorname{coeff}_{[q_0^{-2}]}
 \prod_{l^e\parallel T}
       \left([1]+[l^2]+\cdots+[l^{2e}]\right)
       \quad\text{in }\mathbb Z[H_m].
\end{equation}
If nonzero, it has a word with $\Omega(q/q_0)\le |H_m|-1$.
\end{proposition}
\begin{proof}
The two divisibilities are exactly $q^2\mid R$ and $d\mid q^2$.
Write $q=q_0v$, with $v\mid T$, for the coefficient formula. The preceding
partial-product proof applies to the squared residue of each original prime.
Equivalently the labelled map is
$e_v\mapsto[v^2\bmod m]$ from
$\mathbb Z^{(\operatorname{Div}(T))}$ to $\mathbb Z[H_m]$.
\end{proof}

For
\[
 m=2^\alpha\prod_{i=1}^{r}l_i^{e_i},
 \qquad \alpha\ge2,
\]
where the $l_i$ are distinct odd primes,
\begin{equation}\label{eq:groupsize}
 |H_m|=2^{\max(0,\alpha-3)}
       \prod_{i=1}^{r}\frac{l_i^{e_i-1}(l_i-1)}2.
\end{equation}
At an odd prime power,
$(x-1)(x+1)\equiv0\pmod{l_i^{e_i}}$ has just the two unit solutions $\pm1$.
At modulus four there are two; at $2^\alpha$, $\alpha\ge3$, there are four,
since one of $v_2(x-1),v_2(x+1)$ is one and the other must be at least
$\alpha-1$. CRT and the square-map kernel give \eqref{eq:groupsize}.
Consequently
\begin{equation}\label{eq:nine}
 H_m=\{1\}\iff m\mid24\iff K(u)\mid6\iff
 \boxed{u\mid36.}
\end{equation}

\begin{theorem}[Uniform one-step descent for nine exact words]\label{thm:safe}
For every prime $p\equiv1\pmod4$, every proper E or M trace source with
\[
 u\in\{1,2,3,4,6,9,12,18,36\}
\]
has an original same-channel integral return at
\begin{equation}\label{eq:auto}
 R'=R/K(d)^2,\qquad a'=(p+R')/4,\qquad u'=u,
 \quad a'<a.
\end{equation}
Its complete square-deletion family has exactly
\begin{equation}\label{eq:mult}
 \tau(T)=\prod_{l^e\parallel R}
 \left(\left\lfloor\frac e2\right\rfloor
       -\left\lceil\frac{v_l(d)}2\right\rceil+1\right)
\end{equation}
 distinct original targets.
\end{theorem}
\begin{proof}
All the actual square factors are units modulo $m$. Equation \eqref{eq:nine}
therefore makes every choice in \eqref{eq:square} available, including $q=q_0$.
Here $\tau(T)$ is the number of positive divisors of $T$, with
$v_l(T)=\lfloor v_l(R)/2\rfloor-\lceil v_l(d)/2\rceil$.
The target count is the number of divisors of $T$. Distinct $q$ give distinct
$R'$ and hence distinct distinguished denominators. Theorem~\ref{thm:all}
proves integrality, positivity and strict descent.
\end{proof}
This theorem does not assert the existence of a proper trace source at every
prime. For hard primes its E part is vacuous for these nine words: the factors
2 and 3 are residues, whereas a relaxed E gate has negative word character.
The M part is nonvacuous. The theorem's full domain is all $p\equiv1\pmod4$.

\section{The complement is an available word only on its exact gate domain}
At a fixed original $a$, the involution $u^*=a^2/u$ always preserves the integer
box. It preserves the relaxed M gate and swaps the rational tails.
For a relaxed E input it preserves the relaxed gate exactly when
\begin{equation}\label{eq:comp}
 K(R)\mid p^2-1.
\end{equation}
For M the exact identity
\[
 4uG_M(u^*)-pG_M(u)=R(4a+p)
\]
shows that the complement has the same gcd with $R$, because $4u$ and $p$ are
units modulo $R$.  Thus a proper M source has a proper M complement and the
two rational tails are exchanged.  For E, $R\mid G_E^2$ means
$4u\equiv-1\pmod{K(R)}$, and $4u^*+1\equiv1-p^2$ there. This proves both
directions of \eqref{eq:comp}.

A complement flag followed by an allowed factor $k$ maps
$a$ to $(p+R/k)/4$ and decreases it exactly when $k>1$.  This strictness is
automatic for a proper M complement, but not for an E complement.  For example,
$(p,a,R,u,c)=(37,16,27,2,E)$ is proper, whereas its complement $u^*=128$
has $G_E(u^*)=513=19\cdot27$, hence $d^*=1$; its modulus is
$4K(128)=64$, and its only allowed factor is $k=1$.
The original input is recovered from the flag and old $a$ by applying the same
involution. A fixed E word $u=a$ is not counted as two new source words.

\section{Composite examples where the square choice and multiplicities matter}
At the hard prime $p=1009$, take $a=321=3\cdot107$, $R=275$, $u=9$, channel M.
The rational tails are $215926/5,6054/5$ and have sum $44396$.
Exhaustion of the nine original divisor words $u\mid321^2$ shows that both full
channels of that old shell are empty. Here $d=5$, $m=12$ and $k=25$ gives
\[
 (a',R',u,h',r',s',\lambda')=(255,11,9,1,3,85,8),
\]
\begin{equation}\label{eq:1009}
 \boxed{\frac4{1009}=\frac1{255}+\frac1{686120}+\frac1{24216}.}
\end{equation}
This is a genuine composite distinguished denominator, not the preceding
prime-shell case.

At $p=1108801$, take
\[
 a=279295=5\cdot83\cdot673,\quad R=8379=19\cdot3^2\cdot7^2,
 \quad u=5,\quad c=M.
\]
Its denominator is $d=3$, its budget modulus is 20 and its entire removal set is
\begin{equation}\label{eq:mixed}
 \boxed{\mathcal K=\{21,441\}.}
\end{equation}
The prime 7 is not needed to remove the missing denominator 3, but its residue
is needed to preserve the original divisor word. The minimal square factor
$q_0=3$ fails, whereas $q=3\cdot7$ succeeds. The non-square deletion 21 also
succeeds. The two target records are
\[
\begin{array}{c|rrrrrr}
 k&a'&R'&h'&r'&s'&\lambda'\\\hline
 21&277300&399&5&1&55460&139\\
 441&277205&19&5&1&55441&2918.
\end{array}
\]
The integer coefficient in \eqref{eq:poly} is two.  Its labelled representatives
are $v=7$ and $v=147$, giving $k=21$ and $k=441$.  Thus
$e_7+e_{147}$ survives in the labelled source while its group-ring image
$2[7]$ vanishes after reduction modulo two.  There is one retained original
word $u=5$ and two distinct labelled deletion witnesses and target states.

At $p=51361$, take $a=14845$, $R=8019=11\cdot3^6$, $u=5$, channel M.
Here $d=27$, $q_0=9$, $t=27$ and $m=20$. The only removal is $k=81$, giving
\[
 (a',R',h',r',s',\lambda')=(12865,99,5,1,2573,26),
\]
\[
 \boxed{\frac4{51361}=\frac1{12865}+\frac1{17179740890}+\frac1{6676930}.}
\]
The returned residual still contains $3^2$. Removing every square instead
would give $R'=11$, $a'=12843$, whose square is not divisible by 5.
Thus maximal square stripping is not a valid substitute for \eqref{eq:square}.

The old composite negative control in \cite{PrimeTrace,Trace} can also be
repaired by a different, explicit map. At $p=1009,a=286,R=135$, its original
E word is $u=7436$, with rational tails $6413/3,168238642/3$.
No same-word removal is available. But $K(R)=45$ divides $p^2-1$, so its
complement $u^*=11$ is another relaxed E source. At that source the complete
removal set is $\{45\}$, not a square deletion. It gives
\[
 (a',R',u,h',r',s',\kappa')=(253,3,11,11,1,23,344),
\]
\[
 \boxed{\frac4{1009}=\frac1{253}+\frac1{87032}+\frac1{3818056}.}
\]
The previously failed map to $R=15,a=256$ remains failed. Its replacement changes
the original word by the recorded complement and then uses a non-square divisor.

\section{The nine-seed list is maximal, even allowing every residual divisor}
We first retain the earlier finite cofactor-capacity estimate.
\begin{lemma}[Capacity obstruction, with original target retained]\label{lem:cap}
Let $H=4j-1>u>0$. If $u\nmid j^2$ and a divisor $W\mid j^2$ satisfies
$W\equiv u\pmod H$, then
\begin{equation}\label{eq:cap}
 H\le u^2-3u+1.
\end{equation}
\end{lemma}
\begin{proof}
Write $W=u+nH$ with $n\ge1$, $V=j^2/W$. The exact equality
$16uV=1+H(H+2-16nV)$ has a positive final factor congruent to one modulo four.
Write it as $4\ell+1$, $\ell\ge0$. Comparison gives
\[
 j=\ell+4nV,\qquad u=(4\ell+1)n+\ell^2/V.
\]
If $\ell=0$, then $j=4uV$, contrary to $u\nmid j^2$. Otherwise
$w=\ell^2/V$ is a positive integer, and $\ell\le(u-2)/4$.
Consequently
\[
 H+1=4\ell+\frac{16\ell^2(u-w)}{(4\ell+1)w}
 \le4\ell+\frac{16\ell^2(u-1)}{4\ell+1}
 \le u^2-3u+2.
\]
The last function is increasing for $\ell\ge0$; substitution of $(u-2)/4$
proves the bound. This is the finite-repair argument of \cite{Capacity},
reproduced with the present original target $u$.
\end{proof}

\begin{theorem}[Exact universal seed classification]\label{thm:optimal}
For a fixed positive integer $u$, the following are equivalent:
\begin{enumerate}
\item $u\mid36$;
\item every proper E or M trace source with word $u$, at every prime
$p\equiv1\pmod4$, admits a same-word residual-divisor return.
\end{enumerate}
For every $u\nmid36$ there are infinitely many hard primes with an available
original M word $u$ and an integral rational first trace, but no such return.
Its middle complement has no such return either. At these examples a specified
complete original middle cofactor box is empty, although the prime has explicit
integer endpoint solutions in both channels.
\end{theorem}
\begin{proof}
The positive implication is Theorem~\ref{thm:safe}. For necessity, put
$m=4K(u)$. By \eqref{eq:nine} there is a unit $b\bmod m$ with $b^2\ne1$.
Replace it by its negative if needed to have $b\equiv1\pmod4$.
Dirichlet's theorem \cite[Theorem 18.1]{Dirichlet} supplies a prime
\[
 \delta\equiv-b^2\pmod m,
 \quad\delta>\max(u,u^2-3u+1,7),
\]
and a distinct prime
\[
 t\equiv b^{-1}\pmod m,\qquad t>\max(u,7).
\]
Then $\delta\equiv3\pmod4$, $t\equiv1\pmod4$, and both primes are coprime
to $L=\operatorname{lcm}(840,m)$. Put $R=\delta t^2$ and choose a prime
$e\equiv3\pmod4$ coprime to $LR$.
The three congruences
\begin{equation}\label{eq:AP}
 p\equiv1\pmod L,\qquad
 p\equiv\delta t-4u\pmod R,\qquad
 p\equiv-1\pmod e
\end{equation}
form one reduced class modulo $LRe$. The middle residue is a unit at both
$\delta$ and $t$, since $u$ is coprime to them. Another application of Dirichlet
therefore supplies infinitely many primes in this class with
$p>\max(R,e,B(u)R)$.

Let $a=(p+R)/4$. Since $R\equiv-b^2b^{-2}=-1\pmod m$, the integer $a$ is
a multiple of $K(u)$, proving the original budget $u\mid a^2$. Also
\[
 \gcd(R,p+4u)=\delta t,\qquad d=t>1.
\]
This is a proper rational trace source at a hard prime, with no full M gate.
The only eligible residual divisors congruent to three modulo four are
$\delta,\delta t,R$. The last fails its gate. For the first two, respectively,
\[
 p+\delta\equiv1-b^2\ne0\pmod m,
 \qquad p+\delta t\equiv1-b\ne0\pmod m.
\]
Their original word budgets fail, so no residual-divisor return exists.

For the middle complement $u^*=a^2/u$, direct valuation comparison gives
\[
 K(u^*)=a/B(u),\qquad 4K(u^*)=(p+R)/B(u)>R.
\]
The complement remains a proper trace source with the same exact denominator.
Indeed, modulo $R$,
\[
 4u(p+4u^*)=4up+16a^2\equiv p(p+4u),
\]
and both $4u$ and $p$ are units.  Hence
$\gcd(R,p+4u^*)=\gcd(R,p+4u)=\delta t$ and $d^*=t>1$.
No divisor $1<k\le R$ can be one modulo $4K(u^*)$, while $k=1$ fails
$d^*\mid k$. Thus choosing the other original rational-tail orientation does
not escape this local obstruction.

The old residual's squarefree part is the prime $\delta$. Put
$j=(\delta+1)/4$. The congruence $\delta+1\not\equiv0\pmod{4K(u)}$ gives
$u\nmid j^2$. A middle cofactor-$\delta$ word would require
$W\mid j^2$ and $W\equiv-p/4\equiv u\pmod\delta$.
The strict size choice of $\delta$ contradicts Lemma~\ref{lem:cap}; thus the
\emph{entire} cofactor box is empty. This is not a claim that the original
residual-$\delta$ shell, which is a different source, is empty.

Finally $e\mid p+1$ yields genuine endpoint states. With
$c=(p+1)/e$ and $j_e=(e+1)/4$, use
\[
 E:\ a_E=u_E=(p+e)/4,\quad(h_E,r_E,s_E,\kappa_E)=(a_E,1,1,c),
\]
\[
 M:\ (h_M,r_M,s_M,\lambda_M)=(j_e,1,c,1),\quad
 a_M=j_ec,\quad R_M=c+1,\quad u_M=j_e.
\]
$p\equiv1\pmod4$ and $e\equiv3\pmod4$ give $c\equiv2\pmod4$, so
$R_M\equiv3\pmod4$; also $R_E=e\equiv3\pmod4$.
The inequalities $p>e$, $c\ge2$, $e\ge3$ give the original ranges. Their
channel-labelled identities follow by multiplication.  In particular the M
tuple $(j_ec,pj_ec,pj_e)$ is not asserted to increase; its increasing
permutation is $(j_ec,pj_e,pj_ec)$.  These are counterexamples to
the proposed descent predicate, not to ES.
\end{proof}

\subsection{A complete hard-prime terminal example}
For $u=8$ choose
\[
 m=16,\quad b=5,\quad\delta=71,\quad t=13,\quad e=11,
 \quad R=11999,\quad L=1680.
\]
The reduced prime progression in \eqref{eq:AP} is
\[
 \boxed{p=48116881+221741520n.}
\]
Its first term is prime, proved by trial division through $6936$. It has
\[
 a=12032220,\quad d=13,\quad u=8,
\]
with rational tails
\[
 \frac{943403747137533690}{13},\qquad
 \frac{627251660716}{13}.
\]
Their sum is $72569567260707262$. At the proposed residual 71 the new
distinguished denominator is $12029238$, whose square is $4\bmod8$.
The original exponent-three requirement at two is lost by exactly one occurrence.
The complement budget modulus is $24064440>11999$ and has no proper deletion.

The complete cofactor-$71$ box is $\operatorname{Div}(18^2)$; its residues are
\[
\begin{array}{c|rrrrrrrr}
 W&1&2&3&4&6&9&12&18\\\hline
 W\bmod71&1&2&3&4&6&9&12&18
\end{array}
\]
\[
\begin{array}{c|rrrrrrr}
 W&27&36&54&81&108&162&324\\\hline
 W\bmod71&27&36&54&10&37&20&40.
\end{array}
\]
There is no target residue eight. The separate endpoint middle state has
$j_e=3,c=4374262$ and denominators $(13122786,p\cdot13122786,3p)$.

\section{What remains and what the finite execution establishes}
The composite-source return is exact; the nine automatic words are a maximal
fixed-seed guarantee for this entire residual-divisor method. Beyond them, the
complete coefficient test \eqref{eq:poly}, its square subfamily, and the trace
complement are useful but not universally occupied. The infinite obstruction
retains an actual original word budget and both middle orientations. Enlarging
the search over the same residual divisors cannot repair those sources.
A new universal arrow must change the word beyond that complement, supply a
non-divisor residual through a proved arithmetic construction, or force an
original coefficient in the complete source. None is assumed here.

The main verifier enumerates every original square-divisor word at all primes
$p\equiv1\pmod4$ through 3000. It checks all proper trace sources, every return
and inverse fibre, actual multiplicities, the unit-square classification through
$u=200$, and the explicit reduced progressions for the 91 nonautomatic seeds
through 100. The latter progression records are not 91 asserted prime values.
The separate implementation imports neither the main program nor predecessor
code: it generates primitive $(h,r,s)$ sources and tests every target residual
divisor directly using rational denominators. It independently reproduces the
complete original source lists and all output records. The command receipts
state the observed outcomes; no independent mathematical review or Lean build
is claimed. No finite scan is used as an all-prime occupancy proof.
